% ╔═══════════════════════════════════════╗
% ║   CHAPITRE 5 : ENSEMBLES ET          ║
% ║   APPLICATIONS                        ║
% ╚═══════════════════════════════════════╝
\chapter{Ensembles et Applications}

\section{Motivations}

Les ensembles sont le langage universel des mathématiques. Tout objet mathématique --- nombre, fonction, espace --- est formalisé comme un ensemble ou une construction sur des ensembles. Ce chapitre pose les fondations sur lesquelles reposent l'algèbre, l'analyse et la cryptographie.

\begin{intuitionbox}[Le paradoxe de Russell]
Au début du \textsc{xx}\textsuperscript{e} siècle, Frege pensait avoir refondé les mathématiques sur des bases logiques. Le jeune Russell lui montra qu'un ensemble contenant <<~tous les ensembles~>> mène à une contradiction. Considérons $F = \{E \in \mathcal{E} \mid E \notin E\}$ : si $F \in F$ alors par définition $F \notin F$ ; si $F \notin F$ alors $F$ vérifie la propriété donc $F \in F$. Contradiction dans les deux cas.

L'énigme équivalente : <<~le barbier rase tous ceux qui ne se rasent pas eux-mêmes. Qui rase le barbier ?~>> --- une telle situation ne peut exister.

Ce paradoxe a conduit à la théorie axiomatique des ensembles (ZFC), qui fournit un cadre rigoureux évitant ces contradictions.
\end{intuitionbox}

\begin{intuitionbox}[Lien avec la cybersécurité]
En cybersécurité, les ensembles structurent tout : un groupe d'utilisateurs est un ensemble, une politique de contrôle d'accès définit des sous-ensembles d'actions autorisées, les relations d'équivalence apparaissent dans l'arithmétique modulaire ($\Z/n\Z$) qui est au c\oe ur de RSA et de la cryptographie à courbes elliptiques. La notion de bijection garantit qu'un chiffrement est inversible (déchiffrable).
\end{intuitionbox}


% ─── SECTION 1 : ENSEMBLES ───
\section{Ensembles}

\subsection{Définir des ensembles}

\begin{defbox}[--- Ensemble]
Un \textbf{ensemble} est une collection d'éléments. On note $x \in E$ si $x$ est un élément de $E$, et $x \notin E$ sinon.

\medskip
L'\textbf{ensemble vide}, noté $\emptyset$, est l'ensemble ne contenant aucun élément.
\end{defbox}

\begin{exbox}[Ensembles fondamentaux]
\begin{itemize}[leftmargin=1.5cm]
  \item $\N = \{0, 1, 2, 3, \ldots\}$ --- entiers naturels.
  \item $\Z = \{\ldots, -2, -1, 0, 1, 2, \ldots\}$ --- entiers relatifs.
  \item $\Q = \left\{\dfrac{p}{q} \;\middle|\; p \in \Z,\; q \in \N \setminus \{0\}\right\}$ --- rationnels.
  \item $\R$ --- réels (contient $\sqrt{2}$, $\pi$, $\ln 2$, etc.).
  \item $\C$ --- nombres complexes.
\end{itemize}
\end{exbox}

On peut aussi définir des ensembles par une propriété :
\[
\{x \in \R \mid |x - 2| < 1\}, \quad \{z \in \C \mid z^5 = 1\}, \quad \{x \in \R \mid 0 \leqslant x \leqslant 1\} = [0,1].
\]


\subsection{Inclusion, union, intersection, complémentaire}

\begin{defbox}[--- Opérations sur les ensembles]
Soient $A, B$ des parties d'un ensemble $E$.

\medskip
\textbf{Inclusion :} $E \subset F$ si tout élément de $E$ est aussi un élément de $F$ :
\[\forall x \in E,\; x \in F.\]

\textbf{Égalité :} $E = F \iff E \subset F \text{ et } F \subset E.$

\medskip
\textbf{Ensemble des parties :} $\mathcal{P}(E)$ est l'ensemble de tous les sous-ensembles de $E$.

\medskip
\textbf{Complémentaire :} $\complement_E A = \{x \in E \mid x \notin A\}$, aussi noté $E \setminus A$.

\medskip
\textbf{Union :} $A \cup B = \{x \in E \mid x \in A \text{ ou } x \in B\}$ (<<~ou~>> inclusif).

\medskip
\textbf{Intersection :} $A \cap B = \{x \in E \mid x \in A \text{ et } x \in B\}$.
\end{defbox}

\begin{exbox}
Si $E = \{1, 2, 3\}$ alors :
\[\mathcal{P}(\{1,2,3\}) = \{\emptyset,\; \{1\},\; \{2\},\; \{3\},\; \{1,2\},\; \{1,3\},\; \{2,3\},\; \{1,2,3\}\}\]
soit $2^3 = 8$ parties.
\end{exbox}


\subsection{Règles de calcul}

\begin{propbox}[--- Règles de calcul ensembliste]
Soient $A, B, C$ des parties d'un ensemble $E$.

\medskip
\textbf{Commutativité :} $A \cap B = B \cap A$, \quad $A \cup B = B \cup A$.

\textbf{Associativité :} $A \cap (B \cap C) = (A \cap B) \cap C$, \quad $A \cup (B \cup C) = (A \cup B) \cup C$.

\textbf{Éléments neutres :} $A \cap E = A$, \quad $A \cup \emptyset = A$, \quad $A \cap \emptyset = \emptyset$.

\textbf{Idempotence :} $A \cap A = A$, \quad $A \cup A = A$.

\textbf{Caractérisation de l'inclusion :} $A \subset B \iff A \cap B = A \iff A \cup B = B$.

\medskip
\textbf{Distributivité :}
\begin{align*}
A \cap (B \cup C) &= (A \cap B) \cup (A \cap C) \\
A \cup (B \cap C) &= (A \cup B) \cap (A \cup C)
\end{align*}

\textbf{Lois de De Morgan :}
\begin{align*}
\complement(A \cap B) &= \complement A \cup \complement B \\
\complement(A \cup B) &= \complement A \cap \complement B
\end{align*}

\textbf{Double complémentaire :} $\complement(\complement A) = A$, \quad et donc $A \subset B \iff \complement B \subset \complement A$.
\end{propbox}

\begin{proofbox}[Preuve de la distributivité $A \cap (B \cup C) = (A \cap B) \cup (A \cap C)$]
$x \in A \cap (B \cup C) \iff x \in A \text{ et } x \in (B \cup C) \iff x \in A \text{ et } (x \in B \text{ ou } x \in C)$.

Par distributivité du <<~et~>> sur le <<~ou~>> (logique, chapitre 1) :

$\iff (x \in A \text{ et } x \in B) \text{ ou } (x \in A \text{ et } x \in C) \iff x \in (A \cap B) \cup (A \cap C)$.
\end{proofbox}

\begin{proofbox}[Preuve de De Morgan $\complement(A \cap B) = \complement A \cup \complement B$]
$x \in \complement(A \cap B) \iff x \notin (A \cap B) \iff \text{non}(x \in A \text{ et } x \in B)$.

Par la loi de De Morgan logique : $\iff \text{non}(x \in A) \text{ ou } \text{non}(x \in B) \iff x \notin A \text{ ou } x \notin B \iff x \in \complement A \cup \complement B$.
\end{proofbox}

\begin{intuitionbox}[Pattern structurel]
Les lois de De Morgan ensemblistes sont \emph{exactement} les lois de De Morgan logiques (chapitre 1), traduites en langage d'ensembles. Chaque preuve ensembliste <<~repasse aux éléments~>> et utilise les connecteurs logiques. C'est le même pattern qui apparaît dans les règles de pare-feu : nier une conjonction de conditions donne une disjonction des négations.
\end{intuitionbox}


\subsection{Produit cartésien}

\begin{defbox}[--- Produit cartésien]
Soient $E$ et $F$ deux ensembles. Le \textbf{produit cartésien} $E \times F$ est l'ensemble des couples $(x, y)$ avec $x \in E$ et $y \in F$ :
\[E \times F = \{(x, y) \mid x \in E,\; y \in F\}.\]
\end{defbox}

\begin{exbox}
$\R^2 = \R \times \R$ est le plan usuel. Le cube unité est $[0,1] \times [0,1] \times [0,1] = \{(x,y,z) \mid 0 \leqslant x,y,z \leqslant 1\}$.
\end{exbox}


% ─── SECTION 2 : APPLICATIONS ───
\section{Applications}

\subsection{Définitions}

\begin{defbox}[--- Application]
Une \textbf{application} (ou fonction) $f : E \to F$ associe à chaque élément $x \in E$ un unique élément $f(x) \in F$.

\medskip
\textbf{Égalité :} $f = g \iff \forall x \in E,\; f(x) = g(x)$.

\medskip
\textbf{Graphe :} $\Gamma_f = \{(x, f(x)) \in E \times F \mid x \in E\}$.
\end{defbox}

\begin{defbox}[--- Composition]
Soient $f : E \to F$ et $g : F \to G$. La \textbf{composée} $g \circ f : E \to G$ est définie par :
\[(g \circ f)(x) = g\bigl(f(x)\bigr).\]
\end{defbox}

\begin{exbox}[Composition concrète]
Soient $f : ]0, +\infty[ \to ]0, +\infty[$ définie par $f(x) = \frac{1}{x}$ et $g : ]0, +\infty[ \to \R$ par $g(x) = \frac{x-1}{x+1}$.

Alors pour tout $x > 0$ :
\[
(g \circ f)(x) = g\!\left(\frac{1}{x}\right) = \frac{\frac{1}{x} - 1}{\frac{1}{x} + 1} = \frac{1 - x}{1 + x} = -g(x).
\]
\end{exbox}

L'\textbf{identité} $\mathrm{id}_E : E \to E$ définie par $x \mapsto x$ est l'application neutre pour la composition. La \textbf{restriction} de $f$ à $A \subset E$ est l'application $f_{|A} : A \to F$ définie par $f_{|A}(x) = f(x)$.


\subsection{Image directe et image réciproque}

\begin{defbox}[--- Image directe et image réciproque]
Soit $f : E \to F$.

\medskip
L'\textbf{image directe} d'une partie $A \subset E$ :
\[f(A) = \{f(x) \mid x \in A\} \subset F.\]

L'\textbf{image réciproque} d'une partie $B \subset F$ :
\[f^{-1}(B) = \{x \in E \mid f(x) \in B\} \subset E.\]
\end{defbox}

\begin{rembox}[Attention]
La notation $f^{-1}(B)$ est un tout : elle existe pour toute application $f$, même non bijective. L'image réciproque d'un singleton $f^{-1}(\{y\})$ peut être vide (si $y$ n'a pas d'antécédent), un singleton, un ensemble à plusieurs éléments, ou même $E$ tout entier (si $f$ est constante de valeur $y$).
\end{rembox}


\subsection{Antécédents}

\begin{defbox}[--- Antécédent]
Soit $y \in F$. Tout élément $x \in E$ tel que $f(x) = y$ est un \textbf{antécédent} de $y$ par $f$. L'ensemble des antécédents de $y$ est $f^{-1}(\{y\})$.
\end{defbox}


% ─── SECTION 3 : INJECTION, SURJECTION, BIJECTION ───
\section{Injection, surjection, bijection}

\subsection{Injection et surjection}

\begin{defbox}[--- Injection]
$f : E \to F$ est \textbf{injective} si :
\[\forall x, x' \in E, \quad f(x) = f(x') \implies x = x'.\]
Autrement dit, tout élément de $F$ a \textbf{au plus un} antécédent.
\end{defbox}

\begin{defbox}[--- Surjection]
$f : E \to F$ est \textbf{surjective} si :
\[\forall y \in F,\; \exists x \in E, \quad y = f(x).\]
Autrement dit, tout élément de $F$ a \textbf{au moins un} antécédent. De façon équivalente : $f(E) = F$.
\end{defbox}

\begin{exbox}[Exemples détaillés]
\textbf{1.} $f_1 : \N \to \Q$ définie par $f_1(x) = \frac{1}{1+x}$.

\emph{Injective :} si $f_1(x) = f_1(x')$ alors $\frac{1}{1+x} = \frac{1}{1+x'}$ donc $1+x = 1+x'$ donc $x = x'$. \checkmark

\emph{Non surjective :} on a toujours $f_1(x) \leqslant 1$, donc $y = 2$ n'a pas d'antécédent.

\medskip
\textbf{2.} $f_2 : \Z \to \N$ définie par $f_2(x) = x^2$.

\emph{Non injective :} $f_2(2) = f_2(-2) = 4$ mais $2 \neq -2$.

\emph{Non surjective :} $y = 3$ n'a pas d'antécédent dans $\Z$ (car $\sqrt{3} \notin \Z$).
\end{exbox}

\begin{intuitionbox}[Reformulation par les antécédents]
C'est la façon la plus intuitive de retenir :
\begin{itemize}[leftmargin=1.5cm]
  \item \textbf{Injective} $\iff$ chaque $y$ a \textbf{au plus 1} antécédent (0 ou 1).
  \item \textbf{Surjective} $\iff$ chaque $y$ a \textbf{au moins 1} antécédent ($\geqslant 1$).
  \item \textbf{Bijective} $\iff$ chaque $y$ a \textbf{exactement 1} antécédent.
\end{itemize}
En crypto : un chiffrement doit être bijectif (chaque texte chiffré a un unique antécédent = le texte clair).
\end{intuitionbox}


\subsection{Bijection}

\begin{defbox}[--- Bijection]
$f : E \to F$ est \textbf{bijective} si elle est à la fois injective et surjective :
\[\forall y \in F,\; \exists!\, x \in E, \quad y = f(x).\]
\end{defbox}

\begin{propbox}[--- Caractérisation par la bijection réciproque]
Soit $f : E \to F$.
\begin{enumerate}[leftmargin=1.5cm]
  \item $f$ est bijective $\iff$ il existe $g : F \to E$ telle que $f \circ g = \mathrm{id}_F$ et $g \circ f = \mathrm{id}_E$.
  \item Si $f$ est bijective, cette application $g$ est unique, elle-même bijective, appelée \textbf{bijection réciproque} de $f$ et notée $f^{-1}$. De plus $(f^{-1})^{-1} = f$.
\end{enumerate}
\end{propbox}

\begin{proofbox}
\textbf{Sens $\Rightarrow$ :} $f$ bijective. Construction de $g$ : pour chaque $y \in F$, la surjectivité donne un $x \in E$ tel que $y = f(x)$ ; on pose $g(y) = x$. On a $f(g(y)) = f(x) = y$ pour tout $y$, donc $f \circ g = \mathrm{id}_F$. On compose à droite par $f$ : $f \circ g \circ f = f$, et l'injectivité de $f$ donne $g \circ f = \mathrm{id}_E$.

\medskip
\textbf{Sens $\Leftarrow$ :} Supposons $g$ existe.
\begin{itemize}[leftmargin=1cm]
  \item \emph{Surjectivité :} soit $y \in F$. Posons $x = g(y)$ ; alors $f(x) = f(g(y)) = \mathrm{id}_F(y) = y$.
  \item \emph{Injectivité :} si $f(x) = f(x')$, on compose par $g$ : $g(f(x)) = g(f(x'))$ donc $\mathrm{id}_E(x) = \mathrm{id}_E(x')$ donc $x = x'$.
\end{itemize}

\medskip
\textbf{Unicité :} si $h$ vérifie aussi $h \circ f = \mathrm{id}_E$ et $f \circ h = \mathrm{id}_F$, alors $f \circ h = f \circ g$, et l'injectivité de $f$ donne $h(y) = g(y)$ pour tout $y$, donc $h = g$.
\end{proofbox}

\begin{exbox}
$f : \R \to ]0, +\infty[$ définie par $f(x) = e^x$ est bijective. Sa réciproque est $f^{-1} = \ln$ :
\[\exp(\ln(y)) = y \;\;\forall y > 0, \qquad \ln(\exp(x)) = x \;\;\forall x \in \R.\]
\end{exbox}

\begin{propbox}[--- Composée de bijections]
Si $f : E \to F$ et $g : F \to G$ sont bijectives, alors $g \circ f$ est bijective et :
\[(g \circ f)^{-1} = f^{-1} \circ g^{-1}.\]
\end{propbox}

\begin{proofbox}
Il existe $u = f^{-1}$ et $v = g^{-1}$ telles que $u \circ f = \mathrm{id}_E$, $f \circ u = \mathrm{id}_F$, $v \circ g = \mathrm{id}_F$, $g \circ v = \mathrm{id}_G$.

Alors $(g \circ f) \circ (u \circ v) = g \circ (f \circ u) \circ v = g \circ \mathrm{id}_F \circ v = g \circ v = \mathrm{id}_G$.

Et $(u \circ v) \circ (g \circ f) = u \circ (v \circ g) \circ f = u \circ \mathrm{id}_F \circ f = u \circ f = \mathrm{id}_E$.

Donc $g \circ f$ est bijective de réciproque $f^{-1} \circ g^{-1}$.
\end{proofbox}

\begin{intuitionbox}[L'ordre s'inverse]
$(g \circ f)^{-1} = f^{-1} \circ g^{-1}$ : pour <<~défaire~>> deux opérations successives, on défait la dernière d'abord. C'est le même principe qu'une pile (LIFO) en informatique, ou que l'inversion d'une chaîne de chiffrement.
\end{intuitionbox}


% ─── SECTION 4 : ENSEMBLES FINIS ───
\section{Ensembles finis}

\subsection{Cardinal}

\begin{defbox}[--- Cardinal]
Un ensemble $E$ est \textbf{fini} s'il existe $n \in \N$ et une bijection de $E$ vers $\{1, 2, \ldots, n\}$. Cet entier $n$, unique, est le \textbf{cardinal} de $E$, noté $\mathrm{Card}\, E$.
\end{defbox}

\begin{propbox}[--- Propriétés du cardinal]
\begin{enumerate}[leftmargin=1.5cm]
  \item $B \subset A$ fini $\implies$ $B$ fini et $\mathrm{Card}\, B \leqslant \mathrm{Card}\, A$.
  \item $A \cap B = \emptyset \implies \mathrm{Card}(A \cup B) = \mathrm{Card}\, A + \mathrm{Card}\, B$.
  \item $B \subset A \implies \mathrm{Card}(A \setminus B) = \mathrm{Card}\, A - \mathrm{Card}\, B$.
  \item \textbf{Formule d'inclusion-exclusion :}
    \[\mathrm{Card}(A \cup B) = \mathrm{Card}\, A + \mathrm{Card}\, B - \mathrm{Card}(A \cap B).\]
\end{enumerate}
\end{propbox}

\begin{proofbox}[Preuve de la formule d'inclusion-exclusion]
On décompose $A \cup B = A \cup (B \setminus (A \cap B))$. Les ensembles $A$ et $B \setminus (A \cap B)$ sont disjoints, donc $\mathrm{Card}(A \cup B) = \mathrm{Card}\, A + \mathrm{Card}(B \setminus (A \cap B)) = \mathrm{Card}\, A + \mathrm{Card}\, B - \mathrm{Card}(A \cap B)$.
\end{proofbox}


\subsection{Injection, surjection, bijection et ensembles finis}

\begin{propbox}[--- Cardinal et types d'application]
Soient $E, F$ deux ensembles finis et $f : E \to F$.
\begin{enumerate}[leftmargin=1.5cm]
  \item $f$ injective $\implies \mathrm{Card}\, E \leqslant \mathrm{Card}\, F$.
  \item $f$ surjective $\implies \mathrm{Card}\, E \geqslant \mathrm{Card}\, F$.
  \item $f$ bijective $\implies \mathrm{Card}\, E = \mathrm{Card}\, F$.
\end{enumerate}
\end{propbox}

\begin{propbox}[--- Équivalence pour ensembles de même cardinal]
Soit $f : E \to F$ avec $\mathrm{Card}\, E = \mathrm{Card}\, F$. Alors :
\[f \text{ injective} \iff f \text{ surjective} \iff f \text{ bijective}.\]
\end{propbox}

\begin{proofbox}
Schéma : on montre (i) $\Rightarrow$ (ii) $\Rightarrow$ (iii) $\Rightarrow$ (i).

\medskip
\textbf{Injective $\Rightarrow$ surjective :} $\mathrm{Card}\, f(E) = \mathrm{Card}\, E = \mathrm{Card}\, F$, et $f(E) \subset F$, donc $f(E) = F$.

\textbf{Surjective $\Rightarrow$ bijective :} par l'absurde, si $f$ non injective, deux éléments distincts ont même image, donc $\mathrm{Card}\, f(E) < \mathrm{Card}\, E$. Or $f(E) = F$, donc $\mathrm{Card}\, F < \mathrm{Card}\, E$. Contradiction.

\textbf{Bijective $\Rightarrow$ injective :} immédiat.
\end{proofbox}

\begin{propbox}[--- Principe des tiroirs (Pigeonhole)]
Si l'on range $n > k$ objets dans $k$ tiroirs, alors au moins un tiroir contient au moins deux objets.
\end{propbox}

\begin{intuitionbox}[Application concrète]
Dans un amphi de 400 étudiants, il y a au moins deux étudiants nés le même jour de l'année ($400 > 366$). Le principe des tiroirs est un outil puissant en combinatoire et en cryptanalyse (par ex.\ le paradoxe des anniversaires pour les collisions de hash).
\end{intuitionbox}


\subsection{Dénombrement d'applications}

Soient $E, F$ des ensembles finis non vides avec $\mathrm{Card}\, E = n$ et $\mathrm{Card}\, F = p$.

\begin{propbox}[--- Nombre d'applications]
Le nombre d'applications de $E$ dans $F$ est $p^n = (\mathrm{Card}\, F)^{\mathrm{Card}\, E}$.
\end{propbox}

\begin{proofbox}[Preuve par récurrence sur $n$]
\textbf{Initialisation :} pour $n = 1$, l'unique élément de $E$ peut être envoyé sur $p$ éléments : $p^1$ applications.

\textbf{Hérédité :} supposons le résultat vrai au rang $n$. Soit $E$ de cardinal $n+1$. On fixe $a \in E$ et on pose $E' = E \setminus \{a\}$ ($n$ éléments). Par hypothèse de récurrence, il y a $p^n$ applications de $E'$ vers $F$. Chacune se prolonge en une application de $E$ vers $F$ par $p$ choix pour l'image de $a$. Total : $p^n \times p = p^{n+1}$.
\end{proofbox}

\begin{propbox}[--- Nombre d'injections]
Le nombre d'injections de $E$ dans $F$ est :
\[p \times (p-1) \times \cdots \times (p - n + 1).\]
\end{propbox}

\begin{proofbox}
Pour $a_1$ : $p$ choix. Pour $a_2$ : $p-1$ choix (image distincte de celle de $a_1$). Pour $a_k$ : $p - (k-1)$ choix. Total : $p(p-1) \cdots (p-n+1)$.
\end{proofbox}

\textbf{Notation factorielle :} $n! = 1 \times 2 \times 3 \times \cdots \times n$, avec $0! = 1$ par convention.

\begin{propbox}[--- Nombre de bijections]
Le nombre de bijections d'un ensemble à $n$ éléments dans lui-même est $n!$
\end{propbox}

\begin{exbox}
Parmi les $5^5 = 3125$ applications de $\{1,2,3,4,5\}$ dans lui-même, seules $5! = 120$ sont bijectives.
\end{exbox}


\subsection{Nombre de sous-ensembles}

\begin{propbox}[--- Cardinal de $\mathcal{P}(E)$]
$\mathrm{Card}\, \mathcal{P}(E) = 2^{\mathrm{Card}\, E} = 2^n$.
\end{propbox}

\begin{proofbox}[Preuve par récurrence sur $n$]
\textbf{Initialisation :} $E = \{a\}$ a 2 parties : $\emptyset$ et $\{a\}$.

\textbf{Hérédité :} soit $E$ de cardinal $n+1$. On fixe $a \in E$. Les parties de $E$ se séparent en deux catégories : celles ne contenant pas $a$ (parties de $E \setminus \{a\}$, au nombre de $2^n$) et celles contenant $a$ (de la forme $\{a\} \cup A'$ avec $A' \subset E \setminus \{a\}$, au nombre de $2^n$). Total : $2^n + 2^n = 2^{n+1}$.
\end{proofbox}


\subsection{Coefficients du binôme de Newton}

\begin{defbox}[--- Coefficient binomial]
Le nombre de parties à $k$ éléments d'un ensemble à $n$ éléments est noté $\displaystyle\binom{n}{k}$.
\end{defbox}

\begin{propbox}[--- Propriétés élémentaires]
\begin{enumerate}[leftmargin=1.5cm]
  \item $\displaystyle\binom{n}{0} = 1$, \quad $\displaystyle\binom{n}{1} = n$, \quad $\displaystyle\binom{n}{n} = 1$.
  \item \textbf{Symétrie :} $\displaystyle\binom{n}{k} = \binom{n}{n-k}$.
  \item \textbf{Somme :} $\displaystyle\sum_{k=0}^{n} \binom{n}{k} = 2^n$.
\end{enumerate}
\end{propbox}

\begin{propbox}[--- Formule de Pascal]
Pour $0 < k < n$ :
\[\binom{n}{k} = \binom{n-1}{k-1} + \binom{n-1}{k}.\]
\end{propbox}

\begin{proofbox}
Soit $E$ de cardinal $n$, $a \in E$, $E' = E \setminus \{a\}$. Les parties à $k$ éléments de $E$ se répartissent en deux types : celles ne contenant pas $a$ (parties à $k$ éléments de $E'$ : il y en a $\binom{n-1}{k}$) et celles contenant $a$ ($\{a\} \cup A'$ avec $A'$ partie à $k-1$ éléments de $E'$ : il y en a $\binom{n-1}{k-1}$). D'où la formule.
\end{proofbox}

Cette formule permet de construire le \textbf{triangle de Pascal} ligne par ligne : chaque coefficient est la somme des deux coefficients au-dessus.

\begin{propbox}[--- Formule explicite]
\[\binom{n}{k} = \frac{n!}{k!\,(n-k)!}\]
\end{propbox}


\subsection{Formule du binôme de Newton}

\begin{thmbox}[--- Binôme de Newton]
Pour $a, b \in \R$ (ou $\C$) et $n \in \N^*$ :
\[(a + b)^n = \sum_{k=0}^{n} \binom{n}{k}\, a^{n-k}\, b^k.\]
\end{thmbox}

\begin{proofbox}[Preuve par récurrence sur $n$]
\textbf{Initialisation :} $(a+b)^1 = \binom{1}{0}a + \binom{1}{1}b$. \checkmark

\textbf{Hérédité :} on suppose la formule vraie au rang $n-1$ et on calcule $(a+b)^n = (a+b)(a+b)^{n-1}$. En développant et regroupant les termes en $a^{n-k}b^k$, le coefficient obtenu est $\binom{n-1}{k} + \binom{n-1}{k-1} = \binom{n}{k}$ par la formule de Pascal.
\end{proofbox}

\begin{exbox}
$(a+b)^2 = a^2 + 2ab + b^2$. \quad $(a+b)^3 = a^3 + 3a^2b + 3ab^2 + b^3$.

En posant $a = 1$, $b = 1$ : $\sum_{k=0}^n \binom{n}{k} = 2^n$. En posant $a = -1$, $b = 1$ : $\sum_{k=0}^n (-1)^k \binom{n}{k} = 0$, ce qui signifie qu'un ensemble à $n$ éléments a autant de parties de cardinal pair que de cardinal impair.
\end{exbox}


% ─── SECTION 5 : RELATION D'ÉQUIVALENCE ───
\section{Relation d'équivalence}

\subsection{Définition}

\begin{defbox}[--- Relation d'équivalence]
Soit $E$ un ensemble et $\mathcal{R}$ une relation sur $E$. C'est une \textbf{relation d'équivalence} si elle est :
\begin{enumerate}[leftmargin=1.5cm]
  \item \textbf{Réflexive :} $\forall x \in E, \; x \mathcal{R} x$.
  \item \textbf{Symétrique :} $\forall x, y \in E, \; x \mathcal{R} y \implies y \mathcal{R} x$.
  \item \textbf{Transitive :} $\forall x, y, z \in E, \; x \mathcal{R} y \text{ et } y \mathcal{R} z \implies x \mathcal{R} z$.
\end{enumerate}
\end{defbox}

\begin{exbox}[Exemples et contre-exemples]
\begin{itemize}[leftmargin=1.5cm]
  \item <<~Être parallèle~>> est une relation d'équivalence sur les droites du plan (réflexive, symétrique, transitive). \checkmark
  \item <<~Être du même âge~>> est une relation d'équivalence. \checkmark
  \item <<~Être perpendiculaire~>> n'en est pas une (pas réflexive, pas transitive). $\times$
  \item La relation $\leqslant$ sur $\R$ n'en est pas une (pas symétrique). $\times$
\end{itemize}
\end{exbox}


\subsection{Classes d'équivalence}

\begin{defbox}[--- Classe d'équivalence]
La \textbf{classe d'équivalence} de $x \in E$ est :
\[\mathrm{cl}(x) = \{y \in E \mid y \mathcal{R} x\}.\]
Si $y \in \mathrm{cl}(x)$, on dit que $y$ est un \textbf{représentant} de la classe.
\end{defbox}

\begin{propbox}[--- Propriétés des classes]
\begin{enumerate}[leftmargin=1.5cm]
  \item $\mathrm{cl}(x) = \mathrm{cl}(y) \iff x \mathcal{R} y$.
  \item Pour tout $x, y \in E$ : soit $\mathrm{cl}(x) = \mathrm{cl}(y)$, soit $\mathrm{cl}(x) \cap \mathrm{cl}(y) = \emptyset$.
  \item Les classes d'équivalence forment une \textbf{partition} de $E$ : $E$ est l'union disjointe de toutes les classes.
\end{enumerate}
\end{propbox}

\begin{rembox}
Une \textbf{partition} de $E$ est une famille $\{E_i\}$ de parties non vides de $E$ telles que $E = \bigcup_i E_i$ et $E_i \cap E_j = \emptyset$ pour $i \neq j$.
\end{rembox}


\subsection{Construction des rationnels}

\begin{exbox}[Les rationnels comme classes d'équivalence]
Sur $E = \Z \times \N^*$, on définit :
\[(p, q) \mathcal{R} (p', q') \iff pq' = p'q.\]

C'est une relation d'équivalence (vérification directe de la réflexivité, symétrie, transitivité). On note $\frac{p}{q} = \mathrm{cl}(p, q)$.

Par exemple $(2, 3) \mathcal{R} (4, 6)$ car $2 \times 6 = 3 \times 4$, donc $\frac{2}{3} = \frac{4}{6}$ : mêmes classes, écritures différentes.

L'ensemble $\Q$ des rationnels est précisément l'ensemble des classes d'équivalence de cette relation.
\end{exbox}

\begin{intuitionbox}[Importance conceptuelle]
Cette construction est fondamentale : on \emph{crée} un nouvel objet mathématique (les fractions) à partir d'objets plus simples (les couples d'entiers) en <<~identifiant~>> ceux qui sont équivalents. Le même procédé est utilisé en algèbre abstraite pour construire les groupes quotients, et en crypto pour construire $\Z/n\Z$.
\end{intuitionbox}


\subsection{L'ensemble $\Z/n\Z$}

Soit $n \geqslant 2$ un entier fixé.

\begin{defbox}[--- Congruence modulo $n$]
On définit sur $\Z$ :
\[a \equiv b \pmod{n} \iff n \mid (a - b)\]
c'est-à-dire $a - b$ est un multiple de $n$.
\end{defbox}

\begin{propbox}[--- La congruence est une relation d'équivalence]
\begin{itemize}[leftmargin=1cm]
  \item \emph{Réflexivité :} $a - a = 0 = 0 \cdot n$, donc $a \equiv a \pmod{n}$.
  \item \emph{Symétrie :} si $a - b = kn$ alors $b - a = (-k)n$.
  \item \emph{Transitivité :} si $a - b = kn$ et $b - c = k'n$ alors $a - c = (k + k')n$.
\end{itemize}
\end{propbox}

La classe d'équivalence de $a$ est :
\[\bar{a} = \mathrm{cl}(a) = \{a + kn \mid k \in \Z\} = a + n\Z.\]

Comme $\overline{n} = \bar{0}$, $\overline{n+1} = \bar{1}$, etc., l'ensemble des classes est :
\[\Z/n\Z = \{\bar{0}, \bar{1}, \bar{2}, \ldots, \overline{n-1}\}\]
qui contient exactement $n$ éléments.

\begin{exbox}[$\Z/7\Z$]
\begin{itemize}[leftmargin=1cm]
  \item $\bar{0} = \{\ldots, -14, -7, 0, 7, 14, 21, \ldots\} = 7\Z$
  \item $\bar{1} = \{\ldots, -13, -6, 1, 8, 15, \ldots\} = 1 + 7\Z$
  \item $\bar{6} = \{\ldots, -8, -1, 6, 13, 20, \ldots\} = 6 + 7\Z$
\end{itemize}
Et $\bar{7} = \bar{0}$, $\bar{10} = \bar{3}$ (car $10 - 3 = 7$), etc.
\end{exbox}

\begin{intuitionbox}[Lien avec le chapitre Arithmétique et la crypto]
$\Z/n\Z$ relie directement ce chapitre à l'arithmétique modulaire étudiée au chapitre 4. Les classes de congruence modulo $n$ sont l'environnement naturel de RSA ($\Z/n\Z$ avec $n = pq$), du petit théorème de Fermat ($a^{p-1} \equiv 1 \pmod{p}$), et du logarithme discret (Diffie-Hellman, ElGamal).

Dans la vie courante : les minutes sont modulo 60, les heures modulo 24, les jours de la semaine modulo 7.
\end{intuitionbox}


% ─── SYNTHÈSE ───
\section{Synthèse personnelle --- Ensembles et Applications}

\begin{intuitionbox}[Connexions identifiées]
\textbf{1. Les ensembles comme langage universel.} Chaque objet mathématique de ce cahier est un ensemble ou vit dans un ensemble : $\N$, $\Z$, $\Q$, $\R$, $\C$. Les opérations (union, intersection, complémentaire) sont les traductions ensemblistes des connecteurs logiques du chapitre 1.

\textbf{2. La bijection comme concept central.} Une bijection garantit une correspondance parfaite entre deux ensembles. En crypto, un algorithme de chiffrement est bijectif (sinon on ne peut pas déchiffrer). Le cardinal permet de quantifier cette correspondance : $\mathrm{Card}\, E = \mathrm{Card}\, F \iff$ il existe une bijection $E \to F$.

\textbf{3. Le dénombrement alimente la cryptographie.} Le nombre de clés possibles ($|\text{espace de clés}|$) détermine la résistance au brute-force. Le coefficient binomial $\binom{n}{k}$ et la factorielle $n!$ sont les outils de base pour évaluer cette taille.

\textbf{4. $\Z/n\Z$ boucle la boucle.} Les relations d'équivalence et les classes de congruence, construites ici à partir de zéro, sont exactement les structures utilisées dans le chapitre 4 (Arithmétique) pour le théorème de Fermat et RSA. Ce chapitre fournit la fondation abstraite, le chapitre 4 l'application concrète.
\end{intuitionbox}


